% !TeX root = ../regression_logistique.tex
\chapter{Notation et vocabulaire}
\label{ann:notation}

\orient{%
retrouver la définition, la dimension et le lieu d'introduction d'un symbole ;
passer du vocabulaire français au vocabulaire anglais de référence}{%
aucun}{%
référence ponctuelle ; revenir au chapitre indiqué dans la dernière colonne}

\section{Indices et dimensions}

\begin{table}[htbp]
\centering
\begin{tabular}{c>{\raggedright\arraybackslash}p{74mm}c}
\toprule
symbole & désignation & introduit \\
\midrule
$m$ & nombre d'observations d'un ensemble ; $1\,600$ pour le fichier entier
  & repère~\ref{ch:probleme} \\
$n$ & nombre de variables explicatives ; $2$ en partie~I, $10$ pour le modèle
    complet & ch.~\ref{ch:nuage} \\
$i$ & indice d'observation, de $1$ à $m$ & ch.~\ref{ch:erreur} \\
$j,k$ & indices de coefficient, de $0$ à $n$ ; l'indice $0$ est celui du biais
  & ch.~\ref{ch:descente} \\
$K$ & nombre de classes ; $4$ ici & ch.~\ref{ch:classifieur} \\
\bottomrule
\end{tabular}
\caption{La somme sur les coefficients commence à $j=0$ parce que le biais est
traité comme un coefficient ordinaire, associé à une colonne de $1$.}
\label{tab:not-indices}
\end{table}

\section{Données}

\begin{table}[htbp]
\centering
\begin{tabular}{c>{\raggedright\arraybackslash}p{74mm}c}
\toprule
symbole & désignation & introduit \\
\midrule
$v$ & note brute, telle qu'elle figure dans le fichier
  & ch.~\ref{ch:standardisation} \\
$\mu$ & moyenne d'une colonne, mesurée sur l'apprentissage
  & ch.~\ref{ch:standardisation} \\
$\sigma$ & écart-type d'une colonne, mesuré sur l'apprentissage
  & ch.~\ref{ch:standardisation} \\
$x_j$ & note standardisée dans la matière $j$, soit $(v-\mu)/\sigma$
  & ch.~\ref{ch:standardisation} \\
$x_0$ & colonne constante égale à $1$, porteuse du biais
  & ch.~\ref{ch:donnees} \\
$\x$ & vecteur $(x_0,x_1,\dots,x_n)$ d'une observation & ch.~\ref{ch:nuage} \\
$x_{ij}$ & note standardisée de l'observation $i$ dans la matière $j$
  & ch.~\ref{ch:descente} \\
$\X$ & tableau des $m$ vecteurs d'observation & ch.~\ref{ch:classifieur} \\
$y_i$ & étiquette de référence de l'observation $i$ : $1$ pour la classe
  traitée, $0$ sinon & ch.~\ref{ch:nuage} \\
\bottomrule
\end{tabular}
\caption{La cote standard $(v-\mu)/\sigma$ est notée $x$ et non $z$, contrairement
à l'usage statistique, la lettre $z$ étant réservée au score du modèle.}
\label{tab:not-donnees}
\end{table}

\section{Paramètres et modèle}

\begin{table}[htbp]
\centering
\begin{tabular}{c>{\raggedright\arraybackslash}p{74mm}c}
\toprule
symbole & désignation & introduit \\
\midrule
$w_0$ & biais, coefficient de la colonne constante & ch.~\ref{ch:nuage} \\
$w_j$ & coefficient de la matière $j$ & ch.~\ref{ch:nuage} \\
$\w$ & vecteur $(w_0,w_1,\dots,w_n)$ de tous les coefficients
  & ch.~\ref{ch:erreur} \\
$w$ & vecteur normal $(w_1,\dots,w_n)$, \emph{biais exclu}, employé dans les
  passages géométriques & ch.~\ref{ch:nuage} \\
$\|w\|$ & norme du seul vecteur normal ; vaut $4{,}89$ pour \Mquatre
  & ch.~\ref{ch:nuage} \\
$\|\w\|$ & norme de tous les coefficients ; vaut $6{,}82$ pour \Mquatre
  & \annexeref{ann:norme} \\
$\alpha$ & taux d'apprentissage & ch.~\ref{ch:descente} \\
$\varepsilon$ & seuil de stagnation du critère d'arrêt & ch.~\ref{ch:entrainer} \\
$s$ & seuil de décision sur la probabilité ; $0{,}5$ ici
  & ch.~\ref{ch:probabilite} \\
$\mu$ (régularisation) & coefficient du terme $\tfrac{\mu}{2}\|\w\|^2$
  & \annexeref{ann:norme} \\
\bottomrule
\end{tabular}
\caption{Deux normes distinctes. $\|w\|$, sans le biais, fixe la raideur de la sigmoïde
le long de la perpendiculaire à la frontière ; $\|\w\|$, biais compris, sert à
la régularisation et à la divergence.}
\label{tab:not-parametres}
\end{table}

\section{Fonctions et quantités calculées}

\begin{table}[htbp]
\centering
\begin{tabular}{c>{\raggedright\arraybackslash}p{74mm}c}
\toprule
symbole & désignation & introduit \\
\midrule
$z$ & score, $\sum_j w_jx_j$ & ch.~\ref{ch:nuage} \\
$\sigma(z)$ & fonction logistique, $1/(1+e^{-z})$ & ch.~\ref{ch:probabilite} \\
$p$ & probabilité annoncée pour la classe traitée, $p=\sigma(z)$
  & ch.~\ref{ch:probabilite} \\
$\ell$ & perte associée à une prédiction & ch.~\ref{ch:erreur} \\
$J(\w)$ & perte moyenne sur un ensemble d'observations & ch.~\ref{ch:erreur} \\
$L(\w)$ & vraisemblance & \annexeref{ann:cout} \\
$\nabla J$ & gradient de $J$, vecteur des $n+1$ dérivées partielles
  & ch.~\ref{ch:descente} \\
$\operatorname{softplus}(z)$ & $\log(1+e^{z})$ & \annexeref{ann:stabilite} \\
$\operatorname{logit}(p)$ & $\log\bigl(p/(1-p)\bigr)$ & ch.~\ref{ch:nuage} \\
$d$ & distance signée d'une observation à la frontière, $z/\|w\|$
  & ch.~\ref{ch:nuage} \\
$s_i$ & $\sigma(z_i)(1-\sigma(z_i))$, facteur de la hessienne
  & \annexeref{ann:convexite} \\
\bottomrule
\end{tabular}
\caption{La lettre $\sigma$ désigne l'écart-type d'une colonne dans les passages
sur la préparation des données, et la fonction logistique partout ailleurs. La
distinction se lit à l'argument : $\sigma$ seul est un écart-type, $\sigma(z)$
est une fonction.}
\label{tab:not-fonctions}
\end{table}

\section{États du modèle}

Le tableau~\ref{tab:etats} du repère~\ref{ch:contrat} nomme les quatre jeux de
coefficients employés : \Mzero, \Mquatre, \Mstop et \Movr. Toute valeur
numérique citée dans le document se rapporte à l'un d'eux, et le nom est rappelé
dans la légende de la figure ou du tableau concerné.

\section{Glossaire}

\begin{table}[htbp]
\centering
\small
\begin{tabular}{>{\raggedright\arraybackslash}p{40mm}>{\raggedright\arraybackslash}p{40mm}>{\raggedright\arraybackslash}p{54mm}}
\toprule
français & anglais & remarque \\
\midrule
cote & odds & rapport $p/(1-p)$, non une probabilité \\
log-cote, logit & log-odds, logit & \\
fonction logistique & logistic function, sigmoid & \\
score & score, logit, linear predictor & « logit » désigne aussi le score dans
  la littérature d'apprentissage \\
entropie croisée binaire & binary cross-entropy & aussi \emph{log-loss} \\
vraisemblance & likelihood & \\
descente de gradient & gradient descent & ici en lot complet, \emph{batch} \\
taux d'apprentissage & learning rate & \\
frontière de décision & decision boundary & \\
vecteur normal & normal vector, weight vector & \\
hyperplan affine & affine hyperplane & \\
un-contre-tous & one-vs-rest, one-vs-all & \cite{rifkin2004} \\
matrice de confusion & confusion matrix & \\
rappel & recall, sensitivity & \\
précision & precision & \\
apprentissage / test & training / test set & \\
fuite de données & data leakage & \\
surapprentissage & overfitting & \\
convexe, strictement convexe & convex, strictly convex & deux propriétés
  distinctes (\annexeref{ann:convexite}) \\
séparabilité linéaire & linear separability & \cite{albert1984} \\
régularisation & regularisation & \\
cote standard & $z$-score, standard score & notée $x$ ici \\
imputation & imputation & \cite{littlerubin2019} \\
\bottomrule
\end{tabular}
\caption{Les sorties console des programmes de la partie~II sont en anglais,
comme les noms de fonctions : c'est la langue des identifiants du projet. Le
texte, lui, emploie systématiquement le terme français.}
\label{tab:glossaire}
\end{table}
